\chapter[Fundamentos Teóricos e Revisão Bibliográfica]{Fundamentos Teóricos e Revisão \mbox{Bibliográfica}}

Esta seção reúne os conceitos e trabalhos que fundamentam o desenvolvimento da interface proposta. A revisão concentra-se em quatro eixos: percepção musical por pessoas surdas ou com deficiência auditiva, uso de estímulos vibrotáteis em contextos musicais, instrumentos musicais digitais acessíveis e dispositivos vestíveis destinados à apresentação de informação tátil. Aplicações em musicoterapia e reabilitação motora são discutidas apenas como contextos correlatos, sem integrar o escopo experimental desta monografia.

\section{Surdez, Deficiência Auditiva e Percepção Musical}

A experiência musical resulta da combinação de características como ritmo, intensidade, timbre, altura e organização temporal. Embora a audição desempenhe papel central nessa experiência, sinais visuais, vibrações e informações proprioceptivas também podem participar da percepção e da interação musical.

Entre pessoas surdas ou com deficiência auditiva, a experiência com música não é homogênea. Grau e tipo de perda auditiva, momento de ocorrência da perda, uso de aparelhos auditivos ou implantes cocleares e experiência musical prévia influenciam as informações disponíveis a cada indivíduo \cite{fletcher2021}. Por essa razão, soluções de acessibilidade não devem pressupor uma única forma de percepção musical.

A percepção de emoção musical por estudantes surdos foi investigada em estudo que identificou contribuição de características como ritmo, textura e timbre para a interpretação da música \cite{darrow2006}. O resultado reforça que aspectos da experiência musical podem permanecer acessíveis mesmo quando a percepção convencional de altura ou melodia é limitada.

Dispositivos de auxílio à audição também apresentam restrições específicas para sinais musicais. A fala e a música possuem requisitos distintos de representação temporal e espectral, e parte das características musicais pode ser transmitida com resolução insuficiente por aparelhos auditivos e implantes cocleares \cite{fletcher2021}. Esse cenário justifica a investigação de modalidades sensoriais complementares, sem assumir que elas substituam a audição.

A proposta desta monografia parte dessa perspectiva. Tato e visão são empregados para representar informações associadas à interação com o instrumento, e não para reproduzir integralmente o conteúdo acústico. A correspondência pretendida é funcional: ocorrência de eventos, intensidade relativa do gesto e organização temporal podem ser apresentadas por canais adicionais enquanto a experiência sonora, quando disponível, permanece independente.

\section{Percepção Vibrotátil da Música}

Estímulos vibrotáteis correspondem a vibrações aplicadas ao corpo de forma intencional para transmitir informação. Essa modalidade tem sido estudada em tecnologias assistivas, interação humano-computador e interfaces musicais porque a pele é capaz de perceber variações mecânicas de intensidade, frequência e duração dentro de limites próprios \cite{paisa2023}.

A organização temporal da música é particularmente compatível com a estimulação tátil. Em crianças com deficiência auditiva, a complementação vibrotátil melhorou a identificação de mudanças rítmicas em comparação com condições sem esse recurso \cite{darrowgoll1989}. Em estudo posterior sobre alterações de altura musical, os resultados apresentaram maior variabilidade, indicando que diferentes propriedades musicais não são transmitidas pelo tato com a mesma eficiência \cite{darrow1992}.

A apresentação de notas musicais diretamente sobre a pele foi avaliada por \textcite{hopkins2016}, que analisou limiares, faixa dinâmica e discriminação de frequência em adultos com audição normal e com deficiência auditiva. O estudo demonstrou que a pele pode receber informação musical por vibração, mas também evidenciou limites que devem ser considerados no projeto de dispositivos hápticos.

Resultados semelhantes aparecem em experimentos sobre timbre. Participantes ouvintes e surdos apresentaram desempenho acima do acaso em determinadas tarefas de discriminação exclusivamente vibrotátil, o que indica que diferenças entre sinais podem produzir padrões táteis distinguíveis \cite{russo2012}. Em outro experimento, participantes surdos e ouvintes sincronizaram movimentos corporais à pulsação de música eletrônica apresentada por uma plataforma vibratória, mostrando que informação temporal pode sustentar sincronização mesmo sem apresentação auditiva \cite{tranchant2017}.

O conjunto desses estudos sustenta o uso do tato como canal complementar, sobretudo para informações temporais e variações de intensidade. Não há, contudo, equivalência geral entre informação sonora e informação vibrotátil. A revisão de \textcite{paisa2023} mostra que desempenho e experiência dependem do atuador, da região corporal, do tipo de sinal, da estratégia de mapeamento e da tarefa proposta.

No sistema desenvolvido, optou-se por não estimar características do áudio. A variável de entrada é obtida diretamente do gesto executado no controlador MIDI. O parâmetro \textit{velocity} é convertido em um comando PWM e passa a representar, no domínio de controle, a intensidade relativa associada ao acionamento do \textit{pad}.

\section{Feedback Háptico e Vibrotátil em Interfaces Musicais}

O termo \textit{feedback} háptico abrange realimentações percebidas pelo tato e pela propriocepção. O \textit{feedback} vibrotátil constitui um caso específico no qual atuadores geram vibrações em contato com o corpo \cite{paisa2023}. Em instrumentos acústicos, respostas mecânicas surgem naturalmente da interação com cordas, superfícies e corpos ressonantes; em instrumentos digitais, parte dessa realimentação pode ser reduzida ou inexistente.

A incorporação de \textit{feedback} háptico a instrumentos digitais tem sido investigada tanto para modificar a sensação de execução quanto para transmitir informação adicional. Diferentes modalidades de realimentação, incluindo força e vibração, alteram a experiência do músico durante tarefas de interação com instrumentos digitais \cite{young2017}.

Em aplicações de acessibilidade, a vibração pode cumprir função informacional. A revisão de 63 trabalhos realizada por \textcite{paisa2023} identificou grande diversidade de atuadores, posições corporais, estratégias de processamento e formas de mapeamento entre sinais auditivos e estímulos táteis. Essa diversidade mostra que a escolha do mapeamento não é um detalhe de implementação, mas uma decisão central do projeto.

A região corporal também interfere na solução. Mãos e dedos apresentam alta relevância perceptiva, mas podem estar ocupados durante a execução musical. Braço e punho permitem apresentar estímulos sem impedir o uso das mãos no instrumento, justificando a escolha de uma estrutura vestível no membro superior para o protótipo desta monografia.

A intensidade do gesto foi escolhida como variável de interesse porque está disponível diretamente no evento MIDI e pode ser associada a diferentes níveis de acionamento sem inferência a partir do áudio. A estratégia não pressupõe que a variação de PWM produza sensação linear de intensidade; essa relação depende da resposta eletromecânica do atuador e da percepção do usuário, aspectos que exigem caracterização específica.

\section{Interfaces Musicais Digitais Acessíveis}

Instrumentos musicais digitais, ou DMIs (\textit{Digital Musical Instruments}), utilizam sensores, interfaces eletrônicas e processamento computacional para relacionar gestos de controle à produção musical. A separação entre gesto, processamento e saída permite configurar formas de interação diferentes das encontradas em instrumentos acústicos.

Quando projetados para ampliar a participação de pessoas com diferentes condições físicas, sensoriais ou cognitivas, esses sistemas são frequentemente descritos como instrumentos musicais digitais acessíveis, ou ADMIs (\textit{Accessible Digital Musical Instruments}). Uma revisão sistemática de 113 publicações identificou 83 instrumentos destinados a práticas inclusivas, com soluções baseadas em controladores tangíveis, sistemas sem contato, interfaces cérebro-computador, rastreamento ocular e dispositivos vestíveis \cite{frid2019}.

A literatura de ADMIs não aponta uma configuração única como adequada a todos os usuários. Personalização, prototipagem iterativa e participação da população-alvo aparecem como elementos recorrentes no desenvolvimento de interfaces inclusivas \cite{frid2019}. Esse princípio é relevante para o presente trabalho porque a etapa realizada é técnica; a adequação perceptiva do mapeamento deverá ser investigada posteriormente com usuários.

O EyeHarp exemplifica a possibilidade de dissociar o gesto convencional do controle musical. A interface emprega rastreamento ocular para permitir execução e composição por pessoas com limitações severas de movimento \cite{vamvakousis2016}. Outras propostas de instrumentos sem uso convencional das mãos discutem acessibilidade não apenas como capacidade de produzir notas, mas também em termos de expressividade e qualidade do controle \cite{davanzo2020}.

A arquitetura desenvolvida nesta monografia segue a mesma lógica de separação funcional, embora mantenha um controlador comercial como entrada. A adaptação ocorre nas saídas: a interação MIDI é acompanhada por uma unidade vestível que fornece respostas vibrotátil e visual relacionadas à intensidade registrada pelo controlador.

\section{Música, Treinamento e Deficiência Auditiva}

A literatura reúne práticas distintas envolvendo música e pessoas surdas ou com deficiência auditiva. Treinamento musical, intervenções baseadas em música e musicoterapia possuem objetivos e protocolos diferentes e, por isso, não devem ser tratados como equivalentes.

Experimentos com complementação vibrotátil mostram que estímulos táteis podem participar de atividades musicais com crianças com deficiência auditiva, especialmente em tarefas rítmicas \cite{darrowgoll1989,darrow1992}. Em crianças usuárias de implante coclear, atividades musicais também têm sido discutidas em relação a desenvolvimento perceptivo e participação musical \cite{gfeller2011}.

Programas de treinamento musical foram avaliados em usuários de aparelhos auditivos e implantes cocleares. Uma revisão sistemática e meta-análise identificou efeitos positivos em algumas medidas de percepção musical, mas também registrou heterogeneidade entre métodos, participantes e protocolos \cite{abshukor2021}. Esses resultados não permitem transferir conclusões diretamente para uma interface vibrotátil específica, mas mostram que a experiência musical mediada por tecnologia constitui campo ativo de investigação.

Nesta monografia, aplicações terapêuticas não são avaliadas. A interface é tratada como plataforma técnica para acessibilidade musical, e qualquer uso em musicoterapia depende de protocolos próprios, participação de profissionais da área e avaliação com a população-alvo.

\section{Música e Reabilitação Motora}

Instrumentos musicais digitais também aparecem em estudos de reabilitação motora. A execução musical envolve movimentos orientados a objetivos, repetição, coordenação e retorno sensorial, características que podem ser incorporadas a atividades de recuperação funcional.

Treinamento com piano MIDI e \textit{pads} eletrônicos foi associado a melhorias em medidas de velocidade, precisão e suavidade de movimento após acidente vascular cerebral \cite{schneider2007}. Outro estudo piloto controlado avaliou terapia apoiada por música na recuperação de funções dos membros superiores e identificou melhora em determinadas medidas motoras \cite{tong2015}. A flexibilidade de interfaces MIDI para adaptação de tarefas e parâmetros também foi discutida no contexto de aplicações médicas \cite{kollar2017}.

Abordagens distintas utilizam a sonificação como forma de realimentação. A associação entre movimentos do braço e respostas sonoras foi investigada em reabilitação neurológica, explorando a relação entre ação motora e retorno sensorial \cite{scholz2016}.

Esses resultados mostram que tecnologias musicais podem sustentar diferentes estratégias de interação. Não constituem, porém, evidência clínica para o protótipo desenvolvido neste trabalho. A reabilitação motora permanece fora da validação realizada e é mencionada apenas como possibilidade de aplicação futura.

\section{Trabalhos Correlatos e Estado da Arte}

Entre os sistemas mais próximos da proposta desta monografia está o Auris System, desenvolvido para oferecer uma representação multissensorial da música a pessoas com deficiência auditiva \cite{araujo2017}. A arquitetura combina áudio processado, estímulos visuais e \textit{feedback} tátil e inclui uma pulseira com múltiplos motores vibratórios. Dados de áudio e MIDI são utilizados para compor a representação musical.

A semelhança principal com o presente trabalho está no uso de um dispositivo corporal e de estímulos táteis associados à música. A diferença está na origem e no uso da variável de controle. No protótipo desenvolvido nesta monografia, a informação é obtida diretamente do gesto no MPD218: a \textit{velocity} é processada no momento do acionamento e utilizada para comandar as saídas física e visual, sem análise do conteúdo acústico.

Interfaces hápticas multimodais e dispositivos vestíveis voltados à acessibilidade musical também foram investigados em processo de co-design com participante surda \cite{cavdir2024}. O estudo destaca que localização do estímulo, movimento corporal e participação do usuário influenciam a experiência e devem integrar etapas posteriores do desenvolvimento.

Em perspectiva mais ampla, \textcite{paisa2023} identificaram diversidade de tecnologias e estratégias de mapeamento em displays táteis destinados à ampliação de sinais auditivos. A revisão também aponta limitações recorrentes, como amostras pequenas e avaliações concentradas em condições laboratoriais. Esses aspectos reforçam a necessidade de distinguir validação técnica de avaliação perceptiva.

A Tabela~\ref{tab:estado_arte} sintetiza os trabalhos que apresentam relação mais direta com as decisões adotadas no projeto.

\begin{table}[htbp]
\centering
\caption{Síntese de trabalhos relacionados à percepção e à interação musical vibrotátil}
\label{tab:estado_arte}
\small
\begin{tabular}{p{2.8cm} p{3.3cm} p{4.1cm} p{5.0cm}}
\hline
\textbf{Trabalho} & \textbf{Público/Contexto} & \textbf{Abordagem} & \textbf{Contribuição para este trabalho} \\
\hline
Darrow e Goll (1989) &
Crianças com deficiência auditiva &
Complementação vibrotátil durante identificação rítmica &
Evidencia a utilidade do tato na apresentação de alterações temporais. \\
\hline
Russo et al. (2012) &
Participantes ouvintes e surdos &
Discriminação de timbres por vibração &
Mostra que diferenças entre sinais podem produzir padrões táteis discrimináveis. \\
\hline
Hopkins et al. (2016) &
Adultos com audição normal e deficiência auditiva &
Apresentação vibrotátil de notas musicais &
Caracteriza capacidades e limites da pele para recepção de informação musical. \\
\hline
Tranchant et al. (2017) &
Participantes surdos e ouvintes &
Música apresentada por plataforma vibratória &
Demonstra sincronização corporal com informação rítmica apresentada pelo tato. \\
\hline
Araujo et al. (2017) &
Pessoas com deficiência auditiva &
Cadeira e pulseira vibrotátil em sistema multissensorial &
Fornece referência de arquitetura corporal voltada à acessibilidade musical. \\
\hline
Frid (2019) &
Práticas musicais inclusivas &
Revisão de instrumentos musicais digitais acessíveis &
Contextualiza a solução no campo dos ADMIs e enfatiza adaptação ao usuário. \\
\hline
Paisa et al. (2023) &
Ampliação tátil de áudio e música &
Revisão de 63 displays táteis &
Sistematiza atuadores, locais de aplicação, mapeamentos e limitações experimentais. \\
\hline
Cavdir (2024) &
Acessibilidade musical &
Interfaces hápticas multimodais e dispositivos vestíveis &
Destaca localização corporal, co-design e experiência multissensorial. \\
\hline
\end{tabular}
\legend{Fonte: Elaboração própria (2026).}
\end{table}

Os trabalhos correlatos demonstram que vibração, estímulos visuais e dispositivos vestíveis já integram diferentes propostas de acessibilidade musical. A contribuição desta monografia não consiste no uso isolado de \textit{feedback} vibrotátil. O foco está na implementação e na validação de uma cadeia específica em que a intensidade do gesto é obtida diretamente de uma interface MIDI, processada em tempo de interação e convertida em um comando comum para saídas tátil e visual.

A unidade vestível incorpora alimentação própria e recebe os comandos por Wi-Fi. Essa configuração mantém separadas a interface de entrada, conectada ao computador, e a saída sensorial posicionada no corpo, permitindo avaliar tecnicamente cada módulo da arquitetura.

\section{Síntese da Literatura e Lacuna de Pesquisa}

A literatura revisada demonstra que o tato pode representar propriedades específicas da atividade musical, sobretudo aquelas relacionadas ao tempo e à intensidade, e que dispositivos vestíveis constituem uma alternativa para apresentar essas informações sem ocupar as mãos do usuário. Os estudos também mostram que desempenho e preferência dependem da estratégia de mapeamento, do atuador, da região corporal e do perfil do participante.

Sistemas multissensoriais voltados à acessibilidade musical já combinam áudio, vibração e recursos visuais. A lacuna tratada nesta monografia é, por isso, tecnológica e experimental: avaliar uma arquitetura em que a intensidade do gesto é adquirida diretamente de um controlador MIDI e utilizada para comandar uma unidade vestível com \textit{feedback} vibrotátil e visual, comunicação sem fio e alimentação local.

A validação realizada não busca demonstrar benefício perceptivo, terapêutico ou clínico. Ela estabelece se a cadeia técnica funciona de maneira coerente e se as variáveis de entrada e saída apresentam respostas mensuráveis nas condições de ensaio. O Capítulo 3 detalha a arquitetura empregada para essa finalidade.
